% Français 9115, batch 1 — Gallica views 1–20.
% Pass: Opus 5 (claude-opus-5), 2026-09-12 — first pass, unchecked against
% the views by a human.
%
% What the twenty views hold. This is the opening of the one digitised volume
% of the « Papiers de Sophie Germain », and against the volume's reputation
% for disorder it opens with a clean, continuous, unbroken piece of work: a
% fair copy, in one regular hand throughout, of the beginning of a treatise
% headed « Exposition des principes qui servent de base à la théorie des
% surfaces élastiques ». Eighteen consecutive pages, no gaps, no loose leaves
% between them, the text running from one view to the next without a break —
% and it stops in mid-sentence at the foot of view 20, which is the batch
% boundary and not the end of the piece. Batch 2 continues it.
%
%   v. 1      The binding. Absent from the output.
%   v. 2      A wrapper leaf, numbered 1, carrying one line in a larger and
%             coarser hand than the memoir's: « Manuscrits de mémoires
%             imprimés ». Not the memoir's hand; an owner's or an archivist's
%             label for what follows. Transcribed.
%   v. 3–9    Title and « Avertissement ». The most personal prose in the
%             batch and the reason to read it: she sets out, in order, that
%             two rival hypotheses are in the field, that one of them has « un
%             nom justement célèbre » behind it, that this is a good reason to
%             distrust her own, and that she has done her best to give hers
%             up. Then the history: Chladni's experiments, the prize, the
%             memoir of January 1814, the Institut's verdict, the memoir of
%             January 1816 and the restriction the Académie attached to its
%             approval, Poisson's memoir of August 1814 and his silence about
%             her work. She announces one correction outright — the sign of
%             the term representing the natural curvature is the opposite of
%             the one she had adopted after Euler (« de sono campanorum ») in
%             the 1816 memoir, so that the observed differences of pitch, far
%             from contradicting her theory, agree with it.
%   v. 9–11   « De la nature des hypothèses mathématiques » — two genres of
%             hypothesis, one a statement of the question itself, the other an
%             arbitrary way of viewing it; only the first carries its evidence
%             into its consequences. Her claim is placed here: hers is of the
%             first kind, the rival's of the second.
%   v. 12–20  « Examen a priori, de la nature des deux hypothèses proposées
%             comme devant servir de bases à la théorie des surfaces
%             élastiques », N° 1 to N° 6 (breaking off inside N° 6).
%             N° 1 d'Alembert's definition of elasticity, and the term
%             −(E−N)δ(E−N). N° 2 the repulsive-force surface, and −I δI.
%             N° 3 the long Euler quotation (« Recherches sur la courbure des
%             surfaces », M. de Berlin 1760 p. 119) on the ambiguity of the
%             curvature of a surface, and the reduction of the sum over all
%             sections to the two principal curvatures. N° 4 the measure of
%             the irregularity of a surface, reduced to ρ − ρ′. N° 5 the
%             identification of her N°s 2 and 4 with the hypothesis of the
%             « savant auteur du mémoire sur les surfaces élastiques », with
%             his own words quoted at length. N° 6 begins the comparison of
%             the two hypotheses and is cut off by the batch.
%
% The numbering on the leaves, and why no \folio{} is written. Every leaf from
% view 3 carries an ink number at the head, in the memoir's own ink, and the
% numbers fall into two series that alternate with the views:
%
%   odd views  3  5  7  9 11 13 15 17 19   →  2  3  4  5  6  7  8  9 10
%   even views 4  6  8 10 12 14 16 18 20   →  2  4  6  8 10 12 14 16 18
%
% and on each of the odd views a second number stands beside the first, struck
% through. The two series are consistent with a leaf count on the rectos and a
% page count on the versos of the same leaves (leaf N, verso = page 2N−2, so
% that view 3 is the first page of leaf 2), the struck numbers on the rectos
% being the odd page numbers cancelled. That reading is this pass's
% reconstruction, not something any mark on the leaf states, and the numbers
% are in ink, not in the BnF's pencil. So no \folio{} is written anywhere in
% this file: the numbers are recorded verbatim in a \note{} at the head of the
% transcription, and a human checking the views should settle whether this
% series is the foliation the literature cites for Français 9115.
%
% Notation. Radii of curvature ρ, ρ′ (principal, of the strained figure),
% R, R′ (principal, of the natural figure), ρ_p, ρ′_p, ρ_q, ρ′_q (of curves cut
% by perpendicular sections); N and E for the natural and the elastic figure
% taken as 1/R + 1/R′ and 1/ρ + 1/ρ′; I for the irregularity of the surface;
% δ for the variation. She writes the second derivatives with a straight d —
% r = d²z/dx², s = d²z/dxdy, t = d²z/dy² — where the modern notation is
% partial; that is set here as she sets it. Angles are θ and A, and « 45 » and
% « 45° » both occur for the same angle. Her own underlining, frequent in the
% quotations, is set \emph{}; \uncertain{} also renders as an underline, so a
% reader of the PDF should take the header's word for which is which.
%
% Nothing here was checked against any published transcription.
\documentclass[11pt,a4paper]{article}
% The preamble shared by every transcription of the mathematicians' manuscripts in Gallica.
%
% It rests on one idea: **uncertainty compounds, it does not resolve.** A
% transcription that smooths over a doubtful word has destroyed the only thing
% it was made for — a reader must be able to tell, at every point, what was on
% the page and what was supplied. The apparatus macros below are therefore the
% critical apparatus, not decoration.
%
% The same commands are recognised by `scripts/render.mjs`, which produces the
% site's reading view, and by `scripts/tei.mjs`. A macro added here must be
% added there too, or rendering fails loudly — which is the wanted behaviour.
%
% Comments here are English, like the rest of the repository. The documents
% this preamble sets are French, and that is deliberate: see the skills.

\ProvidesPackage{germain}[2026/09/15 Transcriptions of mathematicians' manuscripts in Gallica]

\usepackage{amsmath,amssymb,amsthm}
\usepackage{mathrsfs}
% Kept from the parent preamble so the renderer's diagram support stays
% available; these manuscripts draw no commutative diagrams, and nothing here uses it.
\usepackage{tikz-cd}
\usepackage[english,french]{babel}
\usepackage{fontspec}

% --- Greek ------------------------------------------------------------------
% The text face stays fontspec's default, Latin Modern, which has no Greek:
% XeTeX drops every glyph it lacks with a line in the log and nothing on the
% page, and the Greek words the manuscripts quote vanished from the PDFs. So
% the Greek and Greek Extended blocks get a XeTeX character class of their own,
% and entering or leaving a run of it switches the family to CMU Serif — the
% same Computer Modern design, with polytonic Greek — and back. Only the
% family changes: a Greek word in \textit or \textbf keeps its shape and series.
% The transitions to and from every other class carry the tokens that class
% has towards an ordinary letter, so French spacing before « ; » or inside
% « » still holds after a Greek word. No grouping, so no brace or \endgroup
% can come out unbalanced; maths is left alone, since XeTeX inserts nothing
% there. Kept identical in `exercices.sty`.
\makeatletter
\newfontfamily\germainGreekfont{cmun}[
  Extension=.otf, NFSSFamily=germaingreek,
  UprightFont=*rm, ItalicFont=*ti, SlantedFont=*sl,
  BoldFont=*bx, BoldItalicFont=*bi, BoldSlantedFont=*bl]
\newXeTeXintercharclass\germain@greek
\def\germain@greekrange#1#2{%
  \@tempcnta=#1\relax
  \loop
    \XeTeXcharclass\@tempcnta=\germain@greek
    \ifnum\@tempcnta<#2\advance\@tempcnta\@ne\repeat}
\germain@greekrange{"0370}{"03FF}
\germain@greekrange{"1F00}{"1FFF}
\def\germain@greekprev{\rmdefault}
\def\germain@greekin{%
  \edef\germain@greekprev{\f@family}\fontfamily{germaingreek}\selectfont}
\def\germain@greekout{\fontfamily{\germain@greekprev}\selectfont}
\def\germain@greekjoin{%
  \XeTeXinterchartokenstate=\@ne
  \@tempcnta=\z@
  \loop
    \ifnum\@tempcnta=\germain@greek\else
      \XeTeXinterchartoks\germain@greek\@tempcnta=\expandafter{%
        \expandafter\germain@greekout\the\XeTeXinterchartoks\z@\@tempcnta}%
      \XeTeXinterchartoks\@tempcnta\germain@greek=\expandafter{%
        \the\XeTeXinterchartoks\@tempcnta\z@\germain@greekin}%
    \fi
    \ifnum\@tempcnta<4095\advance\@tempcnta\@ne\repeat}
% babel-french sets its own classes, and saves and restores the interchar
% state, each time a language is selected: join again after it does.
\addto\extrasfrench{\germain@greekjoin}
\addto\extrasenglish{\germain@greekjoin}
\AtBeginDocument{\germain@greekjoin}
\makeatother

\usepackage{geometry}
\geometry{a4paper,margin=25mm,marginparwidth=28mm}
\usepackage{marginnote}
\usepackage{xcolor}
\usepackage[normalem]{ulem}
\setlength{\emergencystretch}{3em}
\usepackage{hyperref}
\hypersetup{colorlinks=true,linkcolor=black,urlcolor=black,pdfborder={0 0 0}}

% --- Batch metadata ---------------------------------------------------------
% Taken as they stand by the rendering script, which shows them at the head of
% the reading view. They fix what the file claims to cover. The macro names are
% the parent project's, so that the scripts stay shared; \folder here means the
% volume's slug — `fr-9115` — and \pages counts Gallica views.

\newcommand{\folder}[1]{\gdef\grothFolder{#1}}
\newcommand{\batch}[1]{\gdef\grothBatch{#1}}
\newcommand{\pages}[2]{\gdef\grothFirst{#1}\gdef\grothLast{#2}}
\newcommand{\foldertitle}[1]{\gdef\grothFoldertitle{#1}}

% The holder's own shelfmark, printed as they write it: « Français 9115 ».
\newcommand{\shelfmark}[1]{\gdef\grothShelfmark{#1}}

% The dating, copied verbatim from the holder's catalogue — never inferred
% here. « XIXe siècle » is the BnF's; a pass that dates a leaf from its content
% says so in a \note{}, as its own inference.
\newcommand{\dating}[1]{\gdef\grothDating{#1}}

% The status notice, printed in the title block and repeated as a diagonal
% watermark on every page of the PDF. The manuscripts are in the public
% domain; what this line declares is not a right but a quality — an
% unverified first pass by a machine. The skills write the line; a file
% without it gets no watermark, so leaving it out is visible in the output.
\newcommand{\watermark}[1]{\gdef\grothWatermark{#1}}

\gdef\grothFolder{?}\gdef\grothBatch{}\gdef\grothFirst{?}\gdef\grothLast{?}%
\gdef\grothFoldertitle{}\gdef\grothShelfmark{}\gdef\grothDating{}\gdef\grothWatermark{}

% --- The résumé -------------------------------------------------------------
\newenvironment{resume}
  {\small\setlength{\parskip}{0.45em}}
  {\par\medskip\hrule\medskip}

% --- Keywords ---------------------------------------------------------------
% In English, closing the modernised reading's résumé: the modern vocabulary
% under which to search for what the leaves construct. The manifest extracts
% them to tag the volume — this line is the single source of the tags.
\newcommand{\keywords}[1]{\par\smallskip\noindent
  {\footnotesize\textsc{Keywords} --- \textit{#1}}\par}

% --- The critical apparatus -------------------------------------------------

% The Gallica view number, in the margin. This is the anchor that lets a
% disagreement over a formula be settled by looking at *one* image rather than
% a volume of 750. The site uses it to turn the facsimile as the transcription
% is scrolled. A view is one scanned image: usually one side of a leaf.
\newcommand{\page}[1]{%
  \marginnote{\footnotesize\color{gray!70}\textsf{v.\,#1}}%
  \label{page:#1}\ignorespaces}

% The BnF's pencilled foliation, where the leaf carries it and a pass can read
% it: « 198r ». This is what the literature cites — « ff. 198r–208v » — and
% the one line that turns a citation into a view. Recorded, never inferred:
% a folio a pass cannot read is not written.
\newcommand{\folio}[1]{%
  \marginnote{\footnotesize\color{gray!50}\textsf{f.\,#1}}\ignorespaces}

% The modernised reading's coarser counterpart to \page{}: which views a
% section is drawn from.
\newcommand{\pagerange}[2]{%
  \marginnote{\footnotesize\color{gray!70}\textsf{v.\,#1--#2}}%
  \ignorespaces}

% Illegible: never guessed.
\newcommand{\ill}{\textcolor{red!60!black}{[\,\ldots\,]}}

% A doubtful reading: the word is given, and flagged as uncertain.
\newcommand{\uncertain}[1]{\uline{#1}}

% An editorial addition — a missing letter, an implied word.
\newcommand{\add}[1]{\textcolor{blue!50!black}{[#1]}}

% Struck out by the writer. What was crossed out is part of the document.
\newcommand{\struck}[1]{\sout{#1}}

% The transcriber's note, on the state of the page or the mathematics.
\newcommand{\note}[1]{\par\smallskip{\small\color{gray!60!black}\textit{[#1]}}\par\smallskip}

% A marginal note of hers, distinct from the body of the text.
\newcommand{\marginal}[1]{\par\smallskip\noindent\hspace{1em}%
  {\small\color{gray!60!black}#1}\par\smallskip}

% --- Title ------------------------------------------------------------------
% The title block is French, like the documents themselves.

\AddToHook{shipout/background}{%
  \ifx\grothWatermark\empty\else
    \begin{tikzpicture}[remember picture, overlay]
      \node[rotate=52, scale=3.4, align=center, text=gray!18,
            font=\sffamily\bfseries]
        at (current page.center) {\grothWatermark};
    \end{tikzpicture}%
  \fi}

\AtBeginDocument{%
  \begin{center}
    {\large\bfseries Manuscrits de mathématiciens (Gallica) ---
      \ifx\grothShelfmark\empty\grothFolder\else\grothShelfmark\fi}\\[2pt]
    {\itshape\grothFoldertitle}\\[4pt]
    {\small vues \grothFirst--\grothLast
      \ifx\grothBatch\empty\else\ (lot \grothBatch)\fi}\\[2pt]
    \ifx\grothDating\empty\else
      {\small\color{gray!60!black}Datation du catalogue : \grothDating}\\[2pt]
    \fi
    \ifx\grothWatermark\empty\else
      {\small\bfseries\color{red!45!gray}\grothWatermark}\\[2pt]
    \fi
    {\footnotesize\color{gray!60!black}Transcription établie d'après les images IIIF de Gallica.
     Source gallica.bnf.fr / Bibliothèque nationale de France.\\
     Les lectures incertaines sont soulignées, les illisibles notées [\,\ldots\,],
     les ajouts entre crochets ; \textsf{v.} numérote les vues de Gallica,
     \textsf{f.} la foliotation au crayon de la BnF.}
  \end{center}
  \vspace{1em}}

\endinput


\folder{fr-9115}
\shelfmark{Français 9115}
\batch{1}
\pages{1}{20}
\dating{1801-1900}
\watermark{Lecture automatique\\non vérifiée}
\foldertitle{Papiers de Sophie GERMAIN. — Recueil de dissertations et problèmes mathématiques et physiques.}

\begin{document}

\note{Chaque feuillet, à partir de la vue 3, porte à l'encre un numéro en
tête, et ces numéros se répartissent en deux séries qui alternent avec les
vues : vues 3, 5, 7, 9, 11, 13, 15, 17, 19 → 2, 3, 4, 5, 6, 7, 8, 9, 10 ; vues
4, 6, 8, 10, 12, 14, 16, 18, 20 → 2, 4, 6, 8, 10, 12, 14, 16, 18. Sur chacune
des vues impaires un second nombre, biffé, accompagne le premier. Ces numéros
ne sont pas au crayon : aucune foliotation de la Bibliothèque n'a été lue sur
ces vues, et aucun « f. » n'est donc porté dans cette transcription.}

\page{2}
\note{feuillet de garde, numéroté 1 ; la ligne qui suit est d'une main plus
grande et plus grasse que celle du mémoire}
Manuscrits de mémoires imprimés

\page{3}
\note{cachet « Bibliothèque impériale — Mss » au bas du feuillet}

\section{Exposition des principes qui servent de base à la théorie des surfaces élastiques}

\subsection{Avertissement}

Les phénomènes acoustiques dont on doit la connoissance à Mr Chladny, ont
dirigé l'attention des géomètres vers la question des surfaces élastiques.

Malgré leurs travaux ; il existe encore quelqu'embarras dans le choix des
principes qui doivent servir de base à toute cette théorie.

Deux hypothèses essentiellement différentes ont été proposées. L'une d'elle a
pour appui un nom justement célèbre ; c'est une forte raison de se défier de
celle qui m'appartient : aussi ai-je fait tous mes efforts pour y renoncer.

Je sens que j'ai besoin de m'étayer d'un jugement étranger. Il ne me reste que
ce moyen de dissiper le doute qui me poursuit au milieu des recherches que
j'ai entreprises. En exposant mon sentiment j'y joins mes raisons afin que le
public éclairé les pèse et m'apprenne à en mieux juger.

Si je prens quelquefois le ton affirmatif c'est uniquement pour m'affranchir
de l'expression fatiguante du doute. Il suffit d'avertir une fois le lecteur
que bien loin de prétendre fixer son opinion je sollicite

\page{4}
de sa part l'examen critique de la mienne. On me pardonnera, sans doute, de ne
dissimuler ensuite aucun des avantages que je crois reconnoître dans mon
hypothèse.

Aussi tôt que les premières expériences de Mr Chladny me furent connues, il me
parut que l'analyse pouvoit déterminer les lois auxquelles elles sont
assujetties. J'eus occasion d'apprendre d'un grand géomètre dont les premiers
travaux avoient été consacrés à la théorie du son que cette question contenoit
des difficultés que je n'avois pas même soupçonnées. Je cessai d'y penser.

À l'époque du séjour à Paris de Mr Chladny, la vue de ses expérience excita de
nouveau ma curiosité. J'étudiai le mémoire d'Euler sur le cas linéaire, non
pas certainement dans l'intention de concourir au prix,
\marginal{extraordinaire,} que l'institut proposa alors, mais avec l'unique
désir d'apprécier les difficultés dont les termes mêmes du programme me
renouvelloient l'idée.

Dans le cas linéaire les forces d'élasticité sont hypposées proportionnelles à
la raison inverse du rayon de courbure. L'hypothèse qui donne au lieu de la
raison inverse du rayon de courbure d'une simple courbe ; la somme des raisons
inverses des rayons des deux courbures principales d'une surface me frappa par
son analogie et par sa simplicité.

Je présentai, au concours du mois de janvier 1814

\page{5}
un mémoire qui contient : l'hypothèse que j'avais imaginée ; l'équation de la
plaque vibrante ; et, à l'aide de quelques intégrales particulières, la
comparaison entre les résultats de la théorie et l'expérience.

L'institut a jugé que cette comparaison étoit satisfesante. Me sera-t-il
permis de faire remarquer que la question des surfaces élastiques avoit
dèslors perdu une partie des difficultés mentionnées au programme

\note{appel de note * ; la note occupe le bas des vues 5 et 6 et se lit :
« p. 4 du programme. \emph{Cette différence essentielle entre les questions de
mouvement, considérées sous chacun de ces points de vue, dans le simple cas
linéaire, fait concevoir sur le champ qu'on doit trouver des différences de
même espèce, et surtout une grande augmentation de difficultés lorsqu'on veut
introduire deux dimensions dans le calcul} — et plus bas p. 5 et 6. \emph{Voilà
tout ce que les géomètres ont pu faire sur les problèmes des corps sonores,
considérés dans le cas de deux dimensions, et en introduisant même des
simplifications, qui, on ne peut se le dissimuler, changent l'état naturel des
choses de manière que les résultats de l'analyse n'y peuvent point être
applicables. Ces simplifications hypothétiques sont surtout inadmissibles
lorsqu'il s'agit des surfaces vibrantes métalliques, ou jouissant d'une
élasticité naturelle : prenant le cas le plus simple qui est celui du plan, il
est manifeste qu'on ne peut pas lui appliquer la supposition d'Euler\ldots{}
On n'a donc pas même les équations différentielles du mouvement pour cette
espèce de vibration, en envisageant leurs phénomènes tels que la nature les
donne et la seule recherche de ces équations offroit aux géomètres un sujet de
méditation très intéressant\ldots{}} »}

On étoit en possession d'une \emph{hypothèse qui introduisoit deux dimensions
dans le calcul}. Si cette hypothèse n'étoit pas encore justifiée \emph{a
priori} elle n'étoit au moins appuyée sur aucune \emph{simplification qui
changeat l'état naturel des choses}.

L'ignorance la moins excusable, même chez un simple amateur, m'avoit, à la
vérité, empêchée de déduire, d'une manière régulière mon équation de
l'hypothèse qui lui sert de base. Une pareille faute ne permettoit sans doute
pas à la classe d'accorder une aprobation

\page{6}
entière à mon mémoire ; mais elle ne pouvoit détruire aux yeux des géomètres
la relation qui existe entre l'hypothèse et l'équation. Enfin personne ne
contestoit que l'équation différentielle que j'attribuois au mouvement des
plaques vibrantes n'appartînt réellement à ce genre de vibrations \emph{en
envisageant leurs phénomènes tels que la nature les donne}. La difficulté
sembloit être réduite, ou à traiter mon hypothèse conformément aux règles du
calcul, ou à imaginer une autre hypothèse qui conduisit également à l'équation
que j'avois donnée. On pouvoit encore désirer que l'hypothèse choisie fut
justifiée \emph{a priori}.

Au mois d'août 1814, un membre de la classe lut un \emph{mémoire sur les
surfaces élastiques}. L'auteur adopte une hypothèse nouvelle. Il la traite
avec le talent qui caractérise tous ses ouvrages ; et, en se bornant au cas
des surfaces naturellement planes, il parvient à l'équation générale de ce
genre de surfaces.

Il paroît reconnoître que l'on arriveroit à la même équation en employant la
somme des raisons inverses

\page{7}
des deux rayons de courbures principales.

Cette identité de résultat sembleroit devoir me garantir l'exactitude de la
théorie que j'ai cherché à établir ; mais malheureusement il reste encore
entre la doctrine de ce savant auteur et la mienne, des différences
essentielles qui ne peuvent manquer de jetter le doute dans mon esprit.

Dans le mémoire que j'ai envoyé pour le concours du mois de janvier 1816, j'ai
essayé de justifier \emph{a priori} l'hypothèse que j'avois déja proposée, et
j'ai taché de montrer comment cette hypothèse pouvoit conduire à l'équation
d'une surface élastique de figure quelconque.

L'académie, en accordant le prix à ce mémoire, a pourtant mis une certaine
restriction à son approbation.

Sans avoir jamais bien compris le sens de cette restriction, j'ai dû penser
que je n'avois pas encore réussi à fixer le point de vue général de toute
cette théorie.

D'ailleurs le mémoire sur les surfaces élastiques n'a été publié que depuis
l'époque dont je parle, l'auteur n'y fait aucune mension de mon dernier
travail. J'ai lieu de croire qu'il n'adoptoit pas mon sentiment, et l'opinion
d'un si habile géomètre est devenue pour moi un nouveau motif de défiance.

J'ai longtems attendu qu'en publiant, comme il l'avoit fait esperer, la suite
de ses recherches, ce savant auteur me donnat lieu de reconnoître la source de
mon erreur.

J'entreprens à regret la discussion à laquelle je vais me livrer, elle suppose
de ma part un abandon

\page{8}
d'amour propre que le désir de dissiper mes doutes peut seul expliquer.

Afin de mettre le lecteur à même de prononcer, je me propose d'examiner et de
comparer, sous trois points de vue différens les deux hypothèses proposées.

D'abord \emph{a priori}. Cet examen sera précédé de quelques reflexions
générales sur la nature des hypothèses mathématiques.

En me bornant ensuite au cas des plaques vibrantes j'exposerai les conclusions
analytiques auxquelles conduisent l'une et l'autre hypothèses.

Enfin, je comparerai les résultats du calcul aux phénomènes observés.

Si les principes que je soumets en ce moment au jugement des géomètres, leur
semblent en effet devoir servir de base à la théorie des surfaces élastiques,
jespere être dans la suite en état de publier et l'équation générale de ce
genre de surfaces et une comparaison détaillée de la théorie avec l'expérience
non seulement dans le cas des surfaces planes, où le travail de Mr Chladny m'a
déja aidée d'une manière si puissante, mais encore dans d'autres cas où je me
trouve réduite à mes seules observations.

Je puis annoncer dès à présent que le signe du terme qui représente la
courbure naturelle, est contraire à celui que j'avois adopté, en suivant
l'indication d'Euler (\emph{de sono campanorum}) dans mon mémoire pour 1816,
ensorte que les différences de son observées sur des plaques, avant

\page{9}
et après leur courbure, différences que j'ai trouvées plus sensibles en
employant de plus grandes courbures, loin d'être contraires à la théorie que
je cherchois à établir, comme je l'avois dit N° 15 de ce mémoire, y sont
entièrement conformes.

\subsection{De la nature des hypothèses mathématiques}

On doit distinguer deux genres d'hypothèses. Les unes ne sont autre chose que
l'énoncé même de l'état de la question. L'art de les employer se réduit a
celui de les faire entrer dans le calcul. Elles sont accompagnées de
l'évidence et la transmettent à toutes les conséquences que l'on peut en
tirer.

Les autres résultent d'une manière plus ou moins arbitraire d'envisager la
question. Elles sont assujetties au besoin de l'analyse, il ne s'agit plus
d'une simple traduction du langage ordinaire dans la langue des calculs, il
s'agit d'un choix a faire entre les divers points de vues dont la question est
susceptible et le but de ce choix est de rendre l'analyse applicable.

Il résulte de la nature de pareilles hypothèses que leurs conséquences ont
besoin d'être confirmées par la comparaison avec des données prises ailleurs.

Avant cette comparaison, non seulement on ne sauroit affirmer qu'elles
contiennent tout ce qui est de la

\page{10}
question, mais encore on ne peut pas être sûr qu'elles n'ayent rien introduit
d'étranger ou même de contraire à l'état de la question.

Les mathématiques pures n'admettent que le premier genre d'hypothèses ; aussi
l'équation y est-elle considérée comme une définition complette de l'état de
choses auquel elle s'applique, et n'attend-elle aucune confirmation de la
comparaison à l'\uncertain{observation}.
\note{le dernier mot est écrit sur un grattage}

C'est ainsi, par exemple, que les faits usuels qui résultent de l'emploi du
cercle dans les arts ne sauroient rien ajouter à l'évidence des conséquences
qu'on tire de l'équation de cette courbe.

Les mathématiques mixtes sont souvent obligés d'employer des hypothèses
arbitraires, mais les efforts des géomètres tendent à les en affranchir.

Lorsqu'une théorie physicomathématique est déduite des conditions auxquelles
la question est évidemment assujettie elle tire de cette origine un genre de
certitude qui ne sauroit appartenir aux théories fondées sur des hypothèses
arbitraires. On est en état d'affirmer, avant la comparaison avec l'expérience
que les phénomènes suivront les lois données par le calcul. Il est possible, à
la vérité que les conditions qui sont entrées dans l'analyse ne soient pas les
seules auxquelles la question soit physiquement assujettie ; si l'expérience
révèle

\page{11}
l'existence d'anomalies étrangères aux erreurs inévitables, il restera à
complettir la théorie en y introduisant les conditions qui avoient été omises.
La comparaison avec l'expérience ne sert pas dans ce cas à vérifier si les
suppositions admises appartiennent réellement à la question, mais si elles la
contiennent toute entière.

Les géomètres mettent une telle importance à obtenir le genre de certitude
dont je parle, que même lorsqu'une théorie fondée sur quelque supposition
arbitraire a reçu la sanction de l'expérience, ils croient rendre service à la
science en prouvant que la théorie n'avoit pas besoin du secours d'une telle
hypothèse. C'est ainsi que Mr Biot a prouvé que la théorie des surfaces
tendues est indépendante de la supposition à l'aide de laquelle Euler l'avoit
d'abord établie.

Je dois craindre de me faire illusion en pensant, que mon hypothèse, étrangère
à toute idée arbitraire, peut être regardée comme la définition complette de
l'état de la question des surfaces élastiques et que l'autre hypothèse n'est
applicable a la même question que dans les seuls cas où elle est équivalente a
la mienne. Les \uncertain{§§} suivans seront consacrés au développement de
cette opinion. Je l'abandonnerai sans regret si les géomètres ne dédaignant
pas de l'examiner trouvent qu'elle est mal fondée.

\page{12}

\subsection{Examen a priori, de la nature des deux hypothèses proposées comme devant servir de bases à la théorie des surfaces élastiques}

N° 1 Afin d'éviter les objections qui pourroient tomber sur la définition même
du mot élasticité j'adopterai celle de d'Alembert (encyclopédie méthodique
art. élasticité)

\begin{quote}
« Élasticité ou force élastique en mécanique, propriété ou puissance des corps
naturels au moyen delaquelle ils se rétablissent dans la figure et l'étendue
que quelque cause extérieure leur avoit fait perdre »
\end{quote}

Les raisonnemens suivans nous mèneront à l'hypothèse qui doit renfermer la
mesure des forces d'élasticité.

Il résulte de la définition que les forces d'élasticité tendent à détruire la
différence qui existe entre la figure et l'étendue naturelles des corps doués
de ces forces et celles que les mêmes corps ont été contraints de prendre par
l'action d'une cause extérieure.

Les forces d'élasticité dont un corps quelconque est doué ont donc pour mesure
la différence qui existe entre la figure et l'étendue naturelles de ce corps
et celles qu'une cause extérieure l'a forcé de prendre.

Si nous représentons par $N$ la figure et l'étendue naturelles du corps
élastique et par $E$ la figure et l'étendue qu'une force extérieure l'a forcé
de prendre, l'action des forces d'élasticité sera proportionnelle à la
différence $E-N$. Cette action tend à diminuer la différence $E-N$, le terme
dû aux forces d'élasticité sera donc
\[
  -(E-N)\,\delta(E-N).
\]

\page{13}
N° 2 Imaginons à présent une surface également épaisse dans toute son étendue,
dont tous les points soient animés par des forces répulsives égales
entr'elles, et agissant dans une sphère d'activité d'un rayon donné. L'action
de telles forces tendra à placer chaque molécule à distance égale des
molécules voisines et cette distance aura le rayon donné pour mesure.

La grandeur du rayon donné sera la limite de la force
d'\emph{\uncertain{expansibilité}} d'une telle surface. Quelques soit
\marginal{quelle que} la grandeur du rayon, les forces supposées tendront à
rendre la figure de la surface régulière.

Si la grandeur du rayon est insensible, la force d'expansibilité sera
également insensible ; l'effet de l'action des forces répulsives se réduira à
détruire l'irrégularité de la surface, et sera par conséquent proportionnel à
cette irrégularité. Si nous représentons par $I$ l'irrégularité dont il
s'agit, le terme dû à l'action des forces répulsives sera donc
$-I\,\delta I$.

N° 3. Pour déterminer la valeur analytique des termes $-(N-E)\,\delta(N-E)$ et
$-I\,\delta I$ il est nécessaire de connoître l'expression générale de la
courbure des surfaces.
\note{l'ordre des lettres est inversé par rapport au N° 1, où le même terme
est écrit $-(E-N)\,\delta(E-N)$}

Je partirai de l'état de cette question posée par Euler dans son mémoire
intitulé \emph{Recherches sur la courbure des surfaces} M. de Berlin 1760
p. 119

\begin{quote}
« Pour connoître la courbure des lignes courbes » dit cet illustre géomètre,
« la détermination du rayon
\end{quote}

\page{14}
\begin{quote}
osculateur en fournit la plus juste mesure, en nous présentant pour chaque
point de la courbe un cercle, dont la courbure est précisément la même. Mais,
quand on demande la courbure d'une surface, la question est fort équivoque, et
point du tout susceptible d'une réponse absolue, comme dans le cas précédent.
Il n'y a que les surfaces sphériques dont on puisse mesurer la courbure,
attendu que la courbure d'une sphère est la même que celles de ses grands
cercles, et que son rayon peut en être regardé comme la juste mesure. Mais
pour les autres surfaces on n'en sauroit même comparer la courbure avec celle
d'une sphère comme on peut toujours comparer la courbure d'une ligne courbe
avec celle d'un cercle ; la raison en est évidente puisque, dans chaque point
d'une surface, il peut y avoir une infinité de courbures différentes\ldots{}

« Donc la question sur la courbure des surfaces n'est pas susceptible d'une
réponse simple, mais elle exige à la fois une infinité de déterminations :
car, puisqu'on peut tracer par chaque point d'une surface une infinité de
directions, il faut connoître la courbure selon chacune, avant qu'on puisse se
former une juste idée de la courbure de la surface. Or, par chaque point d'une
surface, on peut faire passer une infinité de sections, et cela non seulement
par rapport à toutes les directions sur la surface même, mais aussi par
rapport a leur inclinaison différente sur la surface\ldots{} »
\end{quote}

\page{15}
Conformément à l'indication d'Euler, je considérerai l'ensemble des courbes
qui résulteroient de la section d'un plan passant par le point donné d'une
surface, et la coupant successivement non seulement dans toutes les
directions, \marginal{mais encore dans chacune de ces directions,} sous toutes
les inclinaisons possibles.

La courbure d'une ligne courbe est exprimée par la raison inverse de son rayon
de courbure ; je chercherai donc \emph{la somme des raisons inverses des
rayons de courbures de toutes les lignes courbes produites par les différentes
sections de la surface}.

Soient $\rho_p$ le rayon de courbure de la courbe produite par une section
perpendiculaire à la surface ; $R$ et $R'$ les rayons de courbures de celles
des courbes qui ayant une tangente commune avec la première sont produites par
les sections inclinées de $\theta$ et $\pi-\theta$ degrés ; on aura
\begin{gather*}
  \frac{1}{R} = \frac{\cos\theta}{\rho_p}, \qquad
  \frac{1}{R'} = \frac{\cos(\pi-\theta)}{\rho_p}, \\
  \frac{1}{R} + \frac{1}{R'} = \frac{\cos\theta + \cos(\pi-\theta)}{\rho_p} = 0 .
\end{gather*}

Par conséquent la somme des raisons inverses des rayons de courbures de toutes
les courbes produites par les sections inclinées sera nulle ; et nous n'aurons
plus à nous occuper que des seules sections perpendiculaires à la surface.

Si $\rho_p$ et $\rho'_p$ sont les rayons de courbures des courbes produites
par deux sections perpendiculaires entr'elles $\rho$ et $\rho'$ les rayons de
courbures principales de la surface, $A$ l'angle que le plan qui contient la
courbe dont le rayon est $\rho_p$ fait avec le plan qui contient la courbe
dont le rayon est $\rho$, nous aurons, en vertu des relations

\page{16}
connues
\begin{align*}
  \frac{1}{\rho_p} &= \frac{1}{\rho}\sin^2 A + \frac{1}{\rho'}\cos^2 A \\
  \frac{1}{\rho'_p} &= \frac{1}{\rho}\sin^2\!\left(\frac{\pi}{2}+A\right)
     + \frac{1}{\rho'}\cos^2\!\left(\frac{\pi}{2}+A\right)
   = \frac{1}{\rho}\cos^2 A + \frac{1}{\rho'}\sin^2 A
\end{align*}
et enfin $\dfrac{1}{\rho_p} + \dfrac{1}{\rho'_p} = \dfrac{1}{\rho} +
\dfrac{1}{\rho'}$, théorème également connu.

Il en résulte que \emph{la somme des raisons inverses des rayons de courbures
de toutes les lignes courbes produites par les différentes sections de la
surface}, se réduit à la somme des raisons inverses des deux rayons de
principales courbures de la même surface prise une infinité de fois. On
sentira aisément que l'idée de l'infini est étrangère à la question, et ne se
présente ici qu'a raison de la répétition d'une seule et même mesure qui est
en effet celle de la courbure de la surface.

Il est facile à présent de trouver l'expression analytique du terme
$-(E-N)\,\delta(E-N)$. Car en nommant $R$ et $R'$ les rayons de courbures
principales de la figure naturelle de la surface, $\rho$ et $\rho'$ les rayons
de courbures principales de la figure que l'action d'une cause extérieure l'a
forcé de prendre, figure qu'on pourroit appeller élastique, parceque ce n'est
que lorsque la surface s'y est pliée qu'elle est douée de forces intérieures
qui lui soient propres ; on aura par ce qui précède
$N = \frac{1}{R} + \frac{1}{R'}$, $E = \frac{1}{\rho} + \frac{1}{\rho'}$, et
par conséquent
\[
  -(E-N)\,\delta(E-N) =
  -\left\{ \frac{1}{\rho} + \frac{1}{\rho'}
    - \left( \frac{1}{R} + \frac{1}{R'} \right) \right\}
  \delta\left\{ \frac{1}{\rho} + \frac{1}{\rho'}
    - \left( \frac{1}{R} + \frac{1}{R'} \right) \right\} .
\]

N° 4 La valeur analytique du terme $-I\,\delta I$ dépend de l'irrégularité de
la surface, et la question de l'irrégularité d'une surface semble encore plus
compliquée que celle

\page{17}
de la courbure. Je crois cependant que l'on peut également obtenir la mesure
de ce genre de quantité.

En considérant toujours la totalité des courbes qui passant par le point donné
sont formées par les diverses sections de la surface, on connoîtra
l'irrégularité de la surface, si on connoît \emph{la différence de courbure de
toutes ces courbes comparées chacune à celle qui est produite par la section
perpendiculaire}.

Par la raison donnée N° 3 on peut exclure les courbes formées par les sections
obliques. En effet la somme des courbures produites par toutes les
inclinaisons possibles du plan passant dans une direction donnée, étant nulle,
la différence entre ces premières courbures et celles également produites par
toutes les inclinaisons possibles du plan passant dans la direction
perpendiculaire, sera nécessairement nulle.

Nous n'aurons donc encore à nous occuper que des seules sections normales.

En faisant comme à l'ordinaire $r = \frac{d^2 z}{dx^2}$,
$s = \frac{d^2 z}{dx\,dy}$, $t = \frac{d^2 z}{dy^2}$ et retenant les
dénominations du N° précédent suivant lesquelles $\rho_p$ et $\rho'_p$ sont
les rayons de courbures des courbes produites par deux sections
perpendiculaires entr'elles $\rho$, $\rho'$ les rayons de principales
courbures de la surface, et $A$ l'angle que le plan qui contient la courbe
dont le rayon est $\rho_p$ fait avec le plan qui contient la courbe dont le
rayon est $\rho$, nous aurons a l'aide des formules connues :
\begin{align*}
  \rho_p &= \frac{\sin^2 A + \cos^2 A}
    {r\cos^2 A + 2s\sin A\cos A + t\sin^2 A} \\
  \rho'_p &= \frac{\sin^2 A + \cos^2 A}
    {r\sin^2 A - 2s\sin A\cos A + t\cos^2 A}
\end{align*}

\page{18}
et par conséquent
\begin{align*}
  \frac{1}{\rho_p} - \frac{1}{\rho'_p}
   &= (r-t)\{\cos^2 A - \sin^2 A\} + 4s\sin A\cos A \\
   &= (r-t)\cos 2A + 2s\sin 2A .
\end{align*}

Soient, à présent, $\rho_q$ et $\rho'_q$ les rayons de courbures de deux
courbes produites par deux sections perpendiculaires entr'elles et
$\frac{\pi}{2} - A$, l'angle que le plan qui contient la courbe dont le rayon
est $\rho_q$ fait avec le plan qui contient la courbe dont le rayon est
$\rho$, nous aurons également
\begin{align*}
  \frac{1}{\rho_q} - \frac{1}{\rho'_q}
   &= (r-t)\cos(\pi-2A) + 2s\sin(\pi-2A) \\
   &= -(r-t)\cos 2A + 2s\sin 2A .
\end{align*}

La somme des deux différences $\frac{1}{\rho_p} - \frac{1}{\rho'_p}$,
$\frac{1}{\rho_q} - \frac{1}{\rho'_q}$ se réduira donc à $4s\sin 2A$.

Si les plans de plus grande et plus petite courbure coïncident avec les plans
des $xz$, $yz$ on a, comme on sait, $s=0$ : ainsi la somme des deux
différences deviendra nulle, et la somme des différences de courbures de
toutes les courbes formées par les diverses sections de la surface, prises en
comparant chaque courbure a celle de la courbe produite par la section
perpendiculaire se réduira enfin à la différence des deux courbures
principales de la surface.

En effet, si on excepte les deux cas de $A = 45°$ et $A = 0$, les plans qui
contiennent les courbes dont les rayons sont $\rho_p$ et $\rho'_p$ ne peuvent
jamais se confondre avec les plans qui contiennent les courbes dont les rayons
sont $\rho_q$ et $\rho'_q$.

\page{19}
Lorsque $A = 45$ ; $\frac{\pi}{2} - A = A$ et, à cause de $s=0$ on a
\[
  \rho_p - \rho'_p = (r-t)\cos 2A = 0 .
\]
Ainsi $\rho - \rho'$ est la seule différence qui ne disparoisse pas ; elle
pourra donc être considérée comme la mesure de l'irrégularité de la surface.
Cela posé, on aura
\[
  -I\,\delta I = -(\rho-\rho')\,\delta(\rho-\rho') .
\]
\note{appel de note * après la première formule ; la note occupe le bas de la
vue et se lit : « Je crois avoir prouvé N° 3 que la courbure d'une surface a
pour mesure la somme des raisons inverses des deux rayons de courbures des
courbes produites par deux quelconques des sections normales perpendiculaires
entr'elles. À cause de $\frac{1}{\rho_p} = \frac{1}{\rho'_p}$,
$2\cdot\frac{1}{\rho}$ sera donc aussi la mesure de la même courbure ;
c'est-à-dire que : quelque soit [en marge : \emph{quelle}] \emph{la courbure
d'une surface elle aura pour mesure la courbure de la sphère dont le rayon est
égal au rayon de la courbe formée par la section normale qui fait un angle de
45° avec les plans de plus grande et de moindre courbures de la même surface.}
Ce théorème m'a paru d'autant plus remarquable qu'il est appuyé sur les
principes qui ont porté Euler a penser (voyez le passage cité N° 3) qu'on ne
sauroit même comparer la courbure d'une surface avec celle d'une sphère. »}

N° 5. L'hypothèse qui fait le sujet des N°s 2 et 4 me paroît être entièrement
semblable a celle qu'a adopté le savant auteur \emph{du mémoire sur les
surfaces élastiques}.

En effet on lit (M. de l'Institut 1812. seconde partie p 192
\uncertain{§§ 11 N° 8} de ce mémoire)

\begin{quote}
\ldots{} « Or quelque soit la cause de cette qualité de la matière
(l'élasticité) elle consiste en une tendance des molécules des corps à se
repousser mutuellement, et l'on peut l'attribuer à une force répulsive qui
s'exerce entre ces points, suivant une certaine fonction de leurs distances.
Il est naturel
\end{quote}

\page{20}
\begin{quote}
de penser que cette force, ainsi que toutes les autres action moléculaires,
n'est sensible que pour des distances imperceptibles ; nous admettrons donc
cette hypothèse, et nous supposerons, en conséquence, que la fonction des
distances qui représente la force élastique, n'a de valeur que pour des
valeurs extrêmement petites de la variable qui exprime les distances, et
qu'elle devient nulle aussitôt que cette variable devient sensible. D'ailleurs
cette action répulsive étant censée venir de tous les points matériels qui
composent la surface, il s'en suit que pour des distances égales, la force
répulsive entre deux points devra être proportionnelle au produit des
épaisseurs de la surface qui leur correspondent ; mais par une raison que l'on
verra bientôt nous nous bornerons à considérer les surfaces également épaisses
dans toute leur étendue. »
\end{quote}

On peut voir N°s 12 et suivans par quelle savante analyse l'auteur, après
avoir pris pour plans des coordonnées ceux des principales courbures de la
surface parvient à introduire la différence
$\frac{1}{\rho} - \frac{1}{\rho'}$, dans ses formules.

C'est a l'aide de la même position des plans coordonnés que nous sommes
arrivés à conclure, N° 4, que l'action des forces d'élasticité sur un des
points de la surface est exprimée, dans la présente hypothèse, par le terme
\[
  -\left( \frac{1}{\rho} - \frac{1}{\rho'} \right)
  \delta\left( \frac{1}{\rho} - \frac{1}{\rho'} \right) .
\]
\note{au N° 4 le même terme est donné sous la forme
$-(\rho-\rho')\,\delta(\rho-\rho')$, sans les raisons inverses}

N° 6 Arrêtons nous un instant à l'examen de la nature des deux hypothèses
proposées.

L'une admet simplement que le corps élastique, et par conséquent aussi chacun
des points matériels dont il est

\note{la vue 20 s'arrête au milieu de la phrase ; le texte se poursuit à la
vue 21, hors de ce lot}

\end{document}
